% !PRINTED 76
\weberpage{76}

2. Since two of the $\vartheta$-squares (12) do not stand in a constant ratio, there is also no equation of the form $F^{(n)}(\Theta_0,\Theta_1) = 0$.

3. The total number of the constants which, by (16), (17), occur in the even and odd functions belonging to one and the same characteristic is exactly equal to the order $m$.

\begin{center}
\textbf{\S\,22. \gesperrt{The Addition Theorem.}}
\end{center}

If $u$, $v$ are two variables, then the products
\[
\vartheta_{g_1,g_2}(u+v)\,\vartheta_{g'_1,g'_2}(u-v)
\]
belong to the $\Theta$-functions of the second order, with the characteristic
\[
(g_1+g'_1,\,g_2+g'_2),
\]
and indeed for each of the two variables $u$, $v$; as functions of $v$ they can therefore be represented linearly by the functions (12) and (15) in \S\,21, so that the variable $u$ occurs in the coefficients. When the characteristic is $(0,0)$, this representation is possible in several ways, since one may choose any two of the functions (12); in the other cases it is possible in only one way. One obtains in this manner 16 formulae, of which, however, 6 are to be derived from the others by mere interchange of $v$ with $-v$, so that only 10 essentially different ones remain. These formulae are derived very easily by first setting them up with indeterminate coefficients and then determining these by putting for $v$ such special values as make one or another of the functions $\vartheta(v)$ vanish. Thus, for example:
\[
\vartheta_{00}(u+v)\,\vartheta_{00}(u-v) = A\,\vartheta_{00}^2(v) + B\,\vartheta_{11}^2(v),
\]
and if one puts $v=0$ and $v=\dfrac{1+\omega}{2}$, then, making use of the formulae of \S\,21, one obtains:
\[
A\,\vartheta_{00}^2 = \vartheta_{00}^2(u), \quad B\,\vartheta_{00}^2 = \vartheta_{11}^2(u),
\]
hence

\weq{1}{\vartheta_{00}^2\,\vartheta_{00}(u+v)\,\vartheta_{00}(u-v) = \vartheta_{00}^2(u)\,\vartheta_{00}^2(v) + \vartheta_{11}^2(u)\,\vartheta_{11}^2(v),}
and if in this one replaces $u$ by
\[
u+\frac{1}{2}, \quad u+\frac{\omega}{2}, \quad u+\frac{1+\omega}{2}
\]
then one forms three further relations, which could also have been derived directly in the same way as (1):
%CONT>>
